\section{Conclusions}
\label{sec:conclusion}

This paper presents the first large-scale empirical characterization of e-scooter riding speed behavior, analyzing 2.78~million GPS trajectories with per-point speed profiles from 52 Korean cities, supplemented by 44.1~million trips over 24 months for longitudinal analysis. We defined a comprehensive suite of safety-relevant speed indicators at the trip, user, spatial, and road infrastructure levels, and applied spatial statistical methods, multi-level regression models, gradient-boosted tree models with SHAP interpretability, and longitudinal GEE analysis to systematically characterize patterns of speeding behavior. We summarize the principal findings across six analytical dimensions.

\textbf{Speeding prevalence and distributional characteristics.} We find that 29.6\% of all trips involve at least one speed reading exceeding the Korean regulatory limit of 25~km/h, but speeding is highly concentrated: 60.9\% of users never speed, while 17.3\% speed on more than half of their trips. Operating mode is the dominant determinant---TUB (turbo) mode accounts for 54\% speeding prevalence compared to below 2\% in STD and ECO modes.

\textbf{Spatial concentration and infrastructure associations.} Speeding exhibits statistically significant spatial clustering in all five analyzed cities (Moran's $I$ from 0.16 to 0.71, all $p < 0.001$). Cycling infrastructure reduces the probability of speeding occurrence (OR $= 0.71$) but paradoxically increases its intensity when it does occur, while the fraction of major roads is consistently protective, likely due to congestion effects.

\textbf{Rider demographics and temporal factors.} Riders under 20 exhibit the highest speeding propensity, with the 30--34 age group showing the strongest protective effect (OR $= 0.70$). After controlling for operating mode, speeding odds are lower during evening and night hours than during midday; raw descriptive patterns showing higher nighttime speeding reflect confounding by TUB mode composition. Provincial variation is substantial (Chungnam 45.8\% vs.\ Ulsan 17.6\%) and persists after controlling for individual-level covariates.

\textbf{Rider typologies.} Gaussian mixture modeling identifies four distinct types: Safe Riders (40\%), Moderate-Risk (26\%), Habitual Speeders (26\%), and Stop-and-Go Riders (8\%). TUB mode usage is overwhelmingly the strongest predictor of Habitual Speeder membership (RRR $= 11.8$).

\textbf{Experience and learning effects.} Longitudinal analysis of 44.1~million trips reveals that speeding increases monotonically with rider experience---from 3.8\% on the first trip to 28.0\% after 500+ trips---driven primarily by progressive adoption of the ungoverned TUB mode (11.5\% at trip~1 to 47.8\% by trip~50) rather than within-mode behavior changes.

\textbf{Trip characteristics, road geometry, and model interpretability.} Speeding prevalence increases strongly with trip distance (13.9\% to 54.2\% across deciles). Service and tertiary roads enable speeding (OR $= 2.74$ and 2.05), while primary/secondary roads are protective (OR $< 0.2$). Higher trajectory curvature reduces speeding prevalence (OR $= 0.716$) but amplifies risk when speeding occurs (5$\times$ higher composite risk score). A LightGBM model with SHAP analysis (AUC $= 0.905$) validates the parametric findings, confirming TUB mode as the dominant predictor (mean SHAP $= 2.33$, 3.6$\times$ the next feature).

Our most consequential finding is the effectiveness of software-based speed governors, established through converging evidence at three levels. The within-subject analysis of 12,046 users demonstrates that the speed governor reduces mean speed by 24\%, maximum speed by 31\%, and virtually eliminates speeding. The difference-in-differences analysis of the December 2023 TUB ban provides population-level causal evidence: each 10 percentage point reduction in TUB availability produces a 5.0 percentage point reduction in city-level speeding (95\% CI: [4.1, 7.4]~pp), with a null placebo test confirming causal validity. Together, these individual-level and population-level results provide robust empirical evidence that mandatory speed governors represent the single most effective intervention for e-scooter speed safety.

We identify four priority policy recommendations. First, \textit{mandatory speed governors}: operators should be required to enforce speed limits through software governors, eliminating ungoverned modes in shared fleets. Second, \textit{targeted geofencing}: the spatial clustering of speeding supports zone-specific speed management, deploying reduced speed limits in identified hotspot areas, with particular attention to both straight segments (high prevalence) and curvy segments (high risk when speeding occurs). Third, \textit{graduated enforcement}: the concentration of speeding among a minority of habitual speeders justifies progressive penalties rather than uniform restrictions. Fourth, \textit{graduated mode access}: the finding that experience-driven speeding is mediated by TUB mode adoption supports restricting new riders to governed modes for their initial trips, with ungoverned access granted only after demonstrating safe riding patterns.

Several directions for future research emerge from this work. Linking speed behavior data to crash and injury outcomes would enable direct estimation of speed-risk relationships for e-scooters, strengthening the evidence base for specific speed limit calibration. Real-time speed monitoring and predictive modeling could support dynamic geofencing systems that adapt speed limits based on current conditions (traffic, weather, events). Extending the analysis to other micromobility modes (e-bikes, mopeds) and other geographic contexts (European and North American cities) would test the generalizability of our findings. Higher-frequency kinematic data (1~Hz or better) would enable more precise characterization of acceleration behavior and near-miss events. Our difference-in-differences analysis provides quasi-experimental causal evidence for the speed governor effect, but randomized mode restriction experiments or regression discontinuity approaches exploiting eligibility thresholds could further strengthen the causal evidence base. Analysis of behavioral substitution---whether riders compensate for the TUB ban through longer trips, different road choices, or riskier acceleration patterns---would test the risk compensation hypothesis and refine the net safety assessment.

More broadly, this study demonstrates that operator telemetry data, analyzed at scale with appropriate statistical methods, can provide the empirical evidence that cities need to move from speculative to evidence-based e-scooter speed regulation. The speed indicator framework, spatial analysis methodology, and modeling approach we develop are fully transferable to other operators and cities, providing a replicable blueprint for data-driven micromobility safety assessment.
